\section{Аналитическая модель затрат на общий контекст}
\label{sec:hybrid-accounting}

\subsection{Условный баланс токенов}
Архитектурная мотивация связана с повторной передачей общего материала. Рассмотрим аддитивную модель учёта, в которой общий блок задачи и контекста вносит $S$ токенов при каждой передаче. Внутренняя конфигурация H передаёт его один раз; внешняя M --- в каждом из $m$ вызовов. Пусть $J_H,J_M$ содержат все остальные входные токены, включая оболочки инструкций и передаваемые ответы, а $O_H,O_M$ --- все выходные токены. Обозначим $R_H=J_H+O_H$, $R_M=J_M+O_M$. Тогда
\begin{equation}
T_H=S+R_H,\qquad T_M=mS+R_M,
\qquad T_M-T_H=(m-1)S+R_M-R_H.
\label{eq:hybrid-balance}
\end{equation}
Это тождество учёта при указанной модели блоков. Оно не предполагает равной длины выхода. В частности, более длинное внутреннее рассуждение может компенсировать сокращение передачи контекста; внешние промежуточные ответы учитываются как при генерации, так и при повторной передаче.

\begin{proposition}[Условное достаточное ресурсное преимущество]
\label{prop:hybrid-bound}
Для $m>1$ предположим, что $\E[R_H-R_M\mid S]\le K$ во всём заданном диапазоне контекста, причём постоянная $K\ge0$ не зависит от $S$ в этом диапазоне. Из~\eqref{eq:hybrid-balance} следует
\[
\E[T_M-T_H\mid S]\ge(m-1)S-K.
\]
Поэтому ожидаемое преимущество внутренней конфигурации по токенам положительно при $S>K/(m-1)$ внутри данного диапазона.
\end{proposition}
\begin{proof}
Вычтем два баланса в~\eqref{eq:hybrid-balance}, возьмём условное ожидание при $S$ и применим предполагаемую верхнюю границу для $\E[R_H-R_M\mid S]$.
\end{proof}

Ограничение остатка является содержательным предположением, а не измеренным свойством модели. Если выход или координационные затраты меняются с контекстом так, что граница нарушается, существование порога не следует. Условие сохраняет мотивацию Hybrid-INoT для длинного общего контекста, явно обозначая её пределы. В текущей выборке длина контекста не варьируется как воздействие и $K$ не оценивается; определить порог преимущества или маршрутизации по ней нельзя. Прежние пилоты с компрессией также не определяют эти величины.

Эмпирический реестр использует фактические счётчики поставщика,
\begin{equation}
T_a=\sum_{j=1}^{n_a}(I_{a,j}+O_{a,j}),
\label{eq:hybrid-observed-ledger}
\end{equation}
а не оценивает суммы по~\eqref{eq:hybrid-balance}. Неизвестный расход остаётся неизвестным. Концептуальный блок $S$ не отождествляется с измеренным аддитивным разбиением входа поставщика: границы токенизации и накладные расходы клиента могут препятствовать такому разбиению. Это позволяет не выдавать теоретический баланс за наблюдаемую декомпозицию. При фиксированной топологии общий член сокращается между ролевым и нейтральным условиями; их ресурсная разность относится к изменениям из-за формулировок и другим остаточным различиям. Именно поэтому нужны оба контроля обозначений.

\subsection{Качество и деньги как отдельные ограничения}
Обозначим проверенный успех через $Q_a=\E[Y_t(a)]$, а ожидаемое число токенов через $\bar T_a=\E[T_a]$. Исходное инженерное понятие токенной полезности можно записать как $U_a=Q_a/\bar T_a$, используя отношение ожиданий на общей полностью наблюдаемой совокупности. Более высокое $U_a$ может сочетаться с худшим $Q_a$ и не подтверждает сохранение качества. Поскольку в исследовании различаются также доступные подвыборки качества и ресурсов, основной отчёт сохраняет отдельно парные разности качества и ресурсные отношения, не строя композитный рейтинг полезности и не назначая непроверенных весов сопровождаемости.

Денежный аналог должен учитывать тариф каждого компонента: входа, кешированного входа и выхода. Оценки $V$ и $V_0$ по API-тарифам в разделе~\ref{sec:heldout-methods} реализуют это требование по сохранённым счётчикам. Преимущество по токенам само по себе не устанавливает денежного преимущества при различающихся долях кеша и выхода. Ни одна из оценок не измеряет счёт подписки или полную стоимость владения; вычисления оценщика, инженерная работа и эксплуатация лежат вне наблюдаемого модельного реестра.

Баланс мотивирует возможный механизм эффективности. Изменяется ли качество, полезны ли обозначения ролей и благоприятствует ли наблюдаемая ресурсная разность внутренней конфигурации --- эмпирические вопросы следующего факторного эксперимента. Аналитическое объяснение не добавляет подтверждающей гипотезы к четырём тестам, зафиксированным до основного запуска.
